\documentclass[11pt]{article}
\usepackage[T1]{fontenc}
\usepackage[utf8]{inputenc}
\usepackage{lmodern}
\usepackage[margin=1in]{geometry}
\usepackage{amsmath,amssymb,amsthm,hyperref}
\newtheorem{theorem}{Theorem}
\newtheorem{lemma}[theorem]{Lemma}
\newcommand{\R}{\mathbb R}
\newcommand{\Z}{\mathbb Z}
\newcommand{\N}{\mathbb N}
\newcommand{\one}{\mathbf 1}
\newcommand{\ee}{\mathrm e}
\newcommand{\Lip}{\operatorname{Lip}}
\newcommand{\poly}{\operatorname{poly}}
\newcommand{\Symb}{\operatorname{Symb}}
\title{A local replacement for the filtration argument in GTZ Lemma 13.2}
\author{Project proof draft: /root/astra\_v14\_baseline\_theta}
\date{7 November 2026}
\begin{document}
\maketitle

\noindent\textbf{Status.} This is a proposed project proof, not an independent
acceptance. It replaces only the multilinearity lemma and its auxiliary
filtration construction. It does not certify the other corrections to the
inverse theorem. All source and consumer bindings are recorded separately.

\section{Statement and conventions}

We use the nilmanifold, polynomial-map and nilcharacter conventions of
\cite[Sections 3, 4 and 6]{GTZ}. In particular all Lie groups below are connected,
simply connected and nilpotent; filtration groups are connected rational Lie
subgroups. A rational subgroup $J$ means that $J\cap\Gamma$ is a lattice in $J$.
The multidegree ordering is coordinatewise. Write
$\one=(1,\ldots,1)\in\N^s$ and $|a|=\sum_j a_j$.
The group $G_a$ is trivial when $a\not\leq\one$.
We use $\ee(t)=\exp(2\pi i t)$ and complex conjugation componentwise in tensor
products. A vector-valued nilcharacter has a unit-vector-valued Lipschitz
test function and a real vertical character, as in \cite[Definition 6.1]{GTZ}.
For limit objects, every defining relation is read in its natural ultrapower
extension, using \cite[Appendix A, Definition A.2 and the following discussion,
and Definition A.3, pp.~1324--1325]{GTZ}. In particular the vertical identity
is required for every limit point $x\in{}^*(G/\Gamma)$ and every
$t\in{}^*G_{\one}$, with the standard homomorphism $\eta$ extended to
$^*\eta$. This is the internal nilcharacter convention used in this proof.
We do not infer this exact identity from an identity only at standard points;
bounded Lipschitz continuity alone would not justify that inference.

\begin{theorem}[The unchanged multilinearity statement]\label{thm:main}
Let $s,m\geq1$ be fixed standard integers. Let $\chi$ be a nilcharacter of
multidegree at most $\one$ on $\Z^s$, or a limit nilcharacter with that
description in the sense of \cite[Sections 5--6]{GTZ}. For fixed homomorphisms
$L_1,L'_1,L_2,\ldots,L_s:\Z^m\longrightarrow\Z$, the three pullbacks below
are nilcharacters of degree at most $s$ and satisfy
\begin{align*}
 [\chi(L_1+L'_1,L_2,\ldots,L_s)]_{\Symb^s}
 ={}&[\chi(L_1,L_2,\ldots,L_s)]_{\Symb^s}\\
 &+[\chi(L'_1,L_2,\ldots,L_s)]_{\Symb^s}.
\end{align*}
The same assertion holds in each coordinate. More explicitly, each component
of
\begin{equation}\label{eq:triple}
 \chi(x+x',y_2,\ldots,y_s)\otimes
 \overline{\chi(x,y_2,\ldots,y_s)}\otimes
 \overline{\chi(x',y_2,\ldots,y_s)}
\end{equation}
is a polynomial nilsequence of degree at most $s-1$ on $\Z^{s+1}$.
The representing nilmanifold is fixed by the original filtered nilmanifold;
its test-function Lipschitz bound is bounded in terms of that fixed data and
the original Lipschitz bound, independently of the polynomial coefficients.
\end{theorem}

The symbol notation abbreviates the degree-$s$ equivalence relation of
\cite[Definitions 6.6 and 6.22]{GTZ}. Thus the explicit assertion about
\eqref{eq:triple} supplies the asserted equality, with no approximation or
loss in degree. It is enough to prove the ordinary statement with uniform
fixed-data bounds and then take limit sequences. We supply that uniformity
below. We do not assume $G_0$ equals the group generated by its positive
multidegree groups.

\section{Subgroups and lattices used in the construction}

We use the following basic nilpotent Lie-group facts recorded in
\cite[Section 3, pp.~1242--1243]{GTZ}: the exponential is a diffeomorphism;
the connected subgroup generated by a set is the exponential of the generated
Lie algebra; and, for a finite set $A$, its discrete generated group is
cocompact in its connected real generated group. The last assertion is applied
only when $A$ lies in an already specified discrete lattice. Here are the
compactness details needed to apply these facts.

\begin{lemma}\label{lem:rational}
Let $M$ be a connected simply connected nilpotent Lie group with lattice
$\Lambda$. Let $J_1,\ldots,J_t$ be connected rational subgroups. Then the
connected subgroup generated by the $J_i$ is closed, simply connected, and
rational in $M$. This also applies to diagonal or coordinate embeddings of
such subgroups into a finite product of $M$.
\end{lemma}
\begin{proof}
First a lattice $\Lambda_J$ in a connected Lie group $J$ is finitely generated.
Choose a proper left-invariant path metric, and $R$ such that every point is
within distance $R$ of some lattice point; compactness of $J/\Lambda_J$,
followed by inversion, gives such an $R$. Follow a path from the identity to a
lattice point in steps of length at most one and choose nearby lattice points,
using the prescribed endpoints. Consecutive chosen lattice points differ by
an element in the compact ball of radius $2R+1$. Its intersection with the
discrete lattice is finite and generates the lattice.

The connected real subgroup generated by this finite generating set is all
of $J$. To justify this last assertion, let $U$ be that subgroup. It contains
$\Lambda_J$, so $J/U$ is compact. A connected subgroup of a simply connected
nilpotent group is closed, and its homogeneous space is a Euclidean space.
One can see the latter directly by a chain of Lie subalgebras from
$\mathfrak u$ to $\mathfrak j$, each of codimension one and an ideal in the
next: the normalizer of a proper subalgebra of a nilpotent Lie algebra strictly
contains it, so such a chain is obtained by choosing one normalizing vector at
each step. Successive exponential coordinates give the Euclidean coordinates
on $J/U$. For completeness, the normalizer assertion follows by taking the
first member of the ascending central series not contained in the subalgebra;
the preceding member is contained in it, so any vector in that member but
outside the subalgebra normalizes it. Hence compactness forces $U=J$.

Choose a finite generating set for each $J_i\cap\Lambda$. Their union $A$
lies in $\Lambda$ and generates all the $J_i$ as a connected real group.
Its abstract generated group is discrete because it lies in $\Lambda$.
The finite-set fact from \cite[p.~1243]{GTZ} therefore makes it a lattice in
$\langle A\rangle_{\R}$. Consequently the larger group
$\Lambda\cap\langle A\rangle_{\R}$ is also a lattice there. A Lie subalgebra
of a nilpotent Lie algebra exponentiates to a closed simply connected subgroup,
which proves the other assertions. For a diagonal or coordinate embedding,
apply the same argument to the images of the finite generating sets, which
lie in the product lattice.
\end{proof}

We will also use the elementary central-quotient consequence explicitly.
If $J$ is such a group, $\Delta$ its lattice, and $C$ a connected central
rational subgroup, then $J/C$ is connected simply connected nilpotent and the
image of $\Delta$ is a lattice. Cocompactness follows from the continuous image
of $J/\Delta$. For discreteness, if lattice cosets tend to the identity coset,
choose local lifts tending to the identity and write the lattice points as
those lifts times elements of $C$. Reduce those elements modulo the lattice
$C\cap\Delta$ into a fixed compact set. The adjusted lattice points now lie
in a compact subset of $J$, so only finitely many occur. Thus their cosets
cannot accumulate at the identity. Simple connectedness follows either from
exponential coordinates or the homotopy sequence with connected kernel $C$.

\section{Bounded normalization at zero}

Write $\chi(n)=F(g(n)\Gamma)$, with the data in
Theorem~\ref{thm:main}. Initially restrict to the fixed subgroup $G_0$ and its
lattice $\Gamma_0=G_0\cap\Gamma$, restricting $F$ through the smooth inclusion
$G_0/\Gamma_0\longrightarrow G/\Gamma$. This only changes a fixed metric
comparison constant. In the following notation $G=G_0$ and $\Gamma=\Gamma_0$.

Choose a fixed compact set of representatives for $G/\Gamma$. Write
$g(0)=q\gamma$, with $q$ in that compact set and $\gamma\in\Gamma$. Set
\begin{equation}\label{eq:normalize}
 h(n)=q^{-1}g(n)\gamma^{-1}
     =q^{-1}\bigl(g(n)g(0)^{-1}\bigr)q,
 \qquad F_q(z\Gamma)=F(qz\Gamma).
\end{equation}
Then $h(0)=1$ and $F_q(h(n)\Gamma)=F(g(n)\Gamma)$ exactly.
Each $G_a$ is normal in $G_0$, because $[G_0,G_a]\subseteq G_a$.
The constant map with value $g(0)^{-1}$ is polynomial into the original
filtration. Product closure of polynomial maps
\cite[Corollary B.4]{GTZ}, followed by the filtered conjugation by $q$, shows
that $h$ is polynomial into the same multidegree filtration. No conjugation
of the lattice is made.

The maps $z\Gamma\mapsto qz\Gamma$ are smooth with uniformly bounded
derivatives for $q$ in the fixed compact set. Thus $F_q$ has a uniform
Lipschitz bound. If $F(tz)=\ee(\eta(t))F(z)$ for $t\in G_{\one}$, then
\begin{equation}\label{eq:vertical}
 F_q(tz)=\ee(\eta_q(t))F_q(z),\qquad
 \eta_q(t)=\eta(qtq^{-1}).
\end{equation}
This is a homomorphism on $G_{\one}$. Its restriction to
$G_{\one}\cap\Gamma$ is integer-valued: put $z=\Gamma$ in
\eqref{eq:vertical} and use both $F_q(t\Gamma)=F_q(\Gamma)$ and the fact that
$F_q(\Gamma)$ is a unit vector. In fact $\eta_q=\eta$. For each fixed
$t\in G_{\one}\cap\Gamma$, the continuous map
$q\mapsto\eta(qtq^{-1})$ on connected $G_0$ has integer values, so it equals
its value $\eta(t)$ at the identity. The lattice generates $G_{\one}$ as a
connected real group, by the argument in Lemma~\ref{lem:rational}.
Two continuous real homomorphisms agreeing on that lattice therefore agree
on $G_{\one}$. This proves the assertion for every $q\in G_0$.
In particular the normalization preserves the same standard vertical
character, including for limit choices of $q$.

\section{The descending auxiliary filtration}

For $a\in\{0,1\}^s$ with $a\ne0$, define the following homomorphisms on
$G_a$ whenever the indicated first coordinate occurs:
\[
 A_a(u)=(u,u,u)\quad(a_1=0),\qquad
 B_a(u)=(u,u,1),\quad C_a(u)=(u,1,u)\quad(a_1=1).
\]
For $1\leq i\leq s$ let $H_i$ be the connected subgroup of $G^3$ generated
by these images for all $a$ with $|a|\geq i$. Set $H_{s+1}=\{1\}$ and
\begin{equation}\label{eq:K}
 K=H_1,\qquad K_0=K_1=H_1,\qquad K_i=H_i\ (i\geq1).
\end{equation}
Every group just defined is fixed by the original filtration, not by $g$.
Lemma~\ref{lem:rational} makes it a connected closed simply connected rational
subgroup of $G^3$. Put $\Delta=K\cap\Gamma^3$.
The sequence $(K_i)$ is descending by its definition.

We check the commutator law. For generators of types $A,A$, $A,B$, and
$A,C$, the coordinatewise commutator has respectively type $A,B,C$ and
belongs to the group with index $a+b$. The weight adds.
If two generators both have first index one, every possibly nontrivial
coordinate commutator belongs to $G_{a+b}=\{1\}$ because
$(a+b)_1=2$. Thus the $B,B$, $B,C$, and $C,C$ commutators are trivial.
These assertions hold for their Lie algebras as well. Bracketing generated
Lie algebras and using the Jacobi identity gives
\[
 [\mathfrak h_i,\mathfrak h_j]\subseteq\mathfrak h_{i+j}
 \quad (i,j\geq1).
\]
Here iterated brackets remain in the indicated space: each bracket is a sum
of iterated brackets of the original generators, and its total original
weight is at least $i+j$. BCH, which terminates for a nilpotent group,
then gives $[H_i,H_j]\subseteq H_{i+j}$. The index-zero assertions also
hold since $K_0=H_1$: for example
$[K_0,K_j]\subseteq H_{j+1}\subseteq H_j=K_j$.
Hence $(K_i)$ is a rational degree filtration, is taut at zero, and has
degree at most $s$.

The top group has a particularly simple exact description:
\begin{equation}\label{eq:top}
 K_s=H_s=\{(uv,u,v):u,v\in G_{\one}\}.
\end{equation}
Indeed $a=\one$ is the only index of weight $s$ and has first coordinate
one. The group $G_{\one}$ is abelian since
$[G_{\one},G_{\one}]\subseteq G_{2\one}=\{1\}$. Thus its $B$ and $C$
images commute and their product is exactly the displayed connected group.
It is central in $K$, also following from the commutator law.

\section{The exact polynomial orbit}

Choose once and for all a total order of $\{0,1\}^s$ extending its partial
order. By \cite[Lemma B.9]{GTZ}, the normalized map has the Taylor expansion
\begin{equation}\label{eq:Taylor}
 h(n)=\prod_{a\in\{0,1\}^s\setminus\{0\}} u_a^{n^a},
 \qquad u_a\in G_a,\qquad n^a=\prod_{j=1}^s n_j^{a_j}.
\end{equation}
The absent constant coefficient equals $h(0)=1$. Define
\[
 P(x,x',y)=
 \bigl(h(x+x',y),h(x,y),h(x',y)\bigr),\qquad
 y=(y_2,\ldots,y_s).
\]
This is exactly the ordered product of the following factors, in the order
of \eqref{eq:Taylor}. If $a_1=0$, its factor is
\[
 A_a(u_a)^{\prod_{j=2}^s y_j^{a_j}}.
\]
If $a_1=1$, its factor is
\[
 B_a(u_a)^{x\prod_{j=2}^s y_j^{a_j}}
 C_a(u_a)^{x'\prod_{j=2}^s y_j^{a_j}}.
\]
The latter identity is exact because the powers of the same element $u_a$
commute. In either case the integer exponent has degree $|a|$ and the fixed
group element lies in $K_{|a|}$. A map $v^{Q(n)}$ of this form is polynomial
into $(K_i)$: its $j$-fold derivative is $v^{\Delta^j Q(n)}$, which lies in
$K_j$ for $j\leq\deg Q$, and is the identity for $j>\deg Q$.
Here derivatives commute within the cyclic one-parameter subgroup generated
by $v$, and $K_{\deg Q}\subseteq K_j$ is used in the first case.
The values also lie in $K_0$. Product closure
\cite[Corollary B.4]{GTZ} now proves that
\begin{equation}\label{eq:poly}
 P\in\poly(\Z^{s+1}_{\N}\longrightarrow K_{\N}).
\end{equation}
This argument proves the assertion for all integer inputs, not just a box.
In particular it verifies that the orbit lies in the normalized ambient
group $K$, which the exact-weight construction at index zero did not ensure.

\section{The lower-degree quotient and its test function}

Let $\overline K=K/K_s$, let $\pi:K\longrightarrow\overline K$ be the
quotient map, and let $\overline\Delta=\pi(\Delta)$.
The central-quotient argument after Lemma~\ref{lem:rational} shows that
$\overline K/\overline\Delta$ is a compact nilmanifold with connected simply
connected nilpotent group. Give it the filtration
$\overline K_i=\pi(K_i)$. It is taut and has degree at most $s-1$.
For each $i\leq s$, $K_s\subseteq K_i$, so
\[
 \overline\Delta\cap\pi(K_i)=\pi(\Delta\cap K_i).
\]
The right side is cocompact in $\pi(K_i)$, by the compact quotient of
$K_i/(\Delta\cap K_i)$; it is discrete as a subgroup of
$\overline\Delta$. Thus the quotient filtration is rational at every index.
Equation~\eqref{eq:poly} remains polynomial after applying $\pi$, directly
from the derivative definition.

Define on $K/\Delta$
\[
 T_q(z_1,z_2,z_3)=
 F_q(z_1\Gamma)\otimes\overline{F_q(z_2\Gamma)}
                    \otimes\overline{F_q(z_3\Gamma)}.
\]
It is well-defined and Lipschitz, by restriction from $(G/\Gamma)^3$.
For $(uv,u,v)\in K_s$, equation~\eqref{eq:vertical} multiplies this tensor
by
\[
 \ee\bigl(\eta_q(uv)-\eta_q(u)-\eta_q(v)\bigr)=1.
\]
Therefore $T_q$ is invariant under $K_s$ and descends to a function
$\overline T_q$ on $\overline K/\overline\Delta$.

For explicit control of its regularity, fix a smooth metric on $K/\Delta$
and average it over the compact central torus $K_s/(K_s\cap\Delta)$.
The averaged metric and the original fixed metric are uniformly equivalent.
The quotient metric makes the quotient a Riemannian submersion. A function
constant on its fibers and Lipschitz upstairs is Lipschitz downstairs with
the same bound for these metrics, by horizontal path lifting. All the
metric comparisons are fixed by the original filtered nilmanifold.
The bound for $T_q$ is at most a fixed multiple of the bound for $F_q$,
since its three factors have norm one. The earlier compact normalization
therefore makes the Lipschitz bound for $\overline T_q$ uniform in $g$.

Finally the expression \eqref{eq:triple} is exactly
\[
 \overline T_q\bigl(\pi(P(x,x',y))\overline\Delta\bigr).
\]
This proves the claimed degree at most $s-1$ representation.

For clarity, the individual factors also have valid degree-$s$
representations. In one copy of $G$, take the connected group $J_i$ generated
by $G_a$ for $|a|\geq i\geq1$, and set $J_0=J_1$.
The same rationality and commutator arguments give a taut rational degree-$s$
filtration on $J_1$, with $J_s=G_{\one}$. The Taylor factors of $h$ and of
each of its linear pullbacks are polynomial for this filtration by the same
integer-polynomial argument. Restricting $F_q$ gives a unit-vector-valued
Lipschitz test function with vertical character $\eta_q$.
Hence those factors are nilcharacters of degree at most $s$.

Substitute the homomorphism
$n\mapsto(L_1(n),L'_1(n),L_2(n),\ldots,L_s(n))$ into the lower-degree
representation. Derivatives, or the explicit polynomial factors above, show
that degree does not increase. The definition of symbol equivalence proves
Theorem~\ref{thm:main}. Relabeling coordinates gives the other coordinate
identities.

\section{Limit data and the actual published consumers}

For a limit nilcharacter the filtered nilmanifold and its metric data are
standard fixed data, and the polynomial and Lipschitz test functions are
limit sequences with bounded Lipschitz norm. Choose the bounded representatives
$q$ separately in each member of those sequences; this gives a limit map
with values in the same compact set. All groups $K_i$, the lattice,
the quotient and the metric constants constructed above are unchanged along
the sequence. The same standard vertical character is retained by the preceding
argument. Representatives for the test functions may be taken to satisfy
their defining unit-norm and vertical-character identities on an
ultrafilter-large set of indices: otherwise, choosing a violating point,
and for the second identity a violating group element, at every bad index
would contradict the exact identity of the limit function at the resulting
limit points. This uses the all-limit quantifiers just specified and the
evaluation rule in \cite[Definition A.3]{GTZ}; it would be invalid with only
standard-point quantifiers. Discard those bad indices. The formulas now hold at every
remaining member, and the test-function bounds are uniform. Their limit is
therefore a limit polynomial nilsequence of
degree at most $s-1$ in precisely the required sense. No compactness argument
selects a nilmanifold depending on an unbounded polynomial coefficient.

The first published use, immediately after Lemma 13.2 on p.~1318, expands
$[\widetilde\chi(n+h,\ldots,n+h)]$ by applying the identity in each
coordinate. Before any use of symmetry the exact expansion is
\[
 [\widetilde\chi(n+h,\ldots,n+h)]
 =\sum_{\omega\in\{0,1\}^s}
 [\widetilde\chi(z_{\omega_1},\ldots,z_{\omega_s})],
 \qquad z_0=n,\ z_1=h.
\]
Every map used here is a fixed integer homomorphism from $\Z^2$.
Thus the coefficient counting and subsequent symmetry step in that argument
receive exactly the same multilinearity statement; no stronger inverse
theorem is asserted by this local repair.

The further explicit appeals in Lemmas 13.3, 13.4 and 13.5 on
pp.~1319--1322 likewise use addition, subtraction and repeated copies of
their variables. These are integer homomorphisms from the displayed finite
Cartesian powers of $\Z$, so the theorem applies without a range or
normalization change. Constant shifts in those arguments still require the
separate symbol-shift assertion in E.8; this proof does not discharge it.
The repaired lemma supplies the multilinearity input to Theorem 7.4,
then to the published inverse-theorem argument. Other corrections, including
the false original Proposition 8.3 and its replacement, remain separate
obligations. Source and current-manuscript consumer spans are recorded in
the accompanying contracts.

\begin{thebibliography}{9}
\bibitem{GTZ}
B. Green, T. Tao and T. Ziegler,
\emph{An inverse theorem for the Gowers $U^{s+1}[N]$-norm},
Annals of Mathematics \textbf{176} (2012), 1231--1372.
The frozen final published PDF is the source for the numbered results used
here. The arXiv v5 correction motivates the auxiliary generator types but is
not used as an accepted theorem in this proof.
\end{thebibliography}
\end{document}
